% Public topic-map snapshot of the instructor's tensor lecture seed.
% Upstream: /Users/setup/Library/CloudStorage/Box-Box/Teaching/6050/LaTeX/misc/tensor.tex
% Upstream SHA-256: d7392a91be71f02b8b21a7c15fb66e14fc1d770b02936894dc2cc23c3053fc61
% Diagrams, executable snippets, device/performance claims, channels-last examples,
% and categorical claims are omitted. This file records only the lecture's topic
% boundary and standard equations; the book prose, code, and figures are independent.
\documentclass{article}
\usepackage{amsmath,amssymb}
\usepackage[margin=1in]{geometry}

\title{Tensors in Computation: Public Topic Map}
\author{Heman Shakeri}
\date{}

\begin{document}
\maketitle

\section*{Public Snapshot Boundary}

The upstream lecture surveys multidimensional arrays, shapes, batched affine layers,
broadcasting, and the distinction between computational arrays and mathematical
tensors. Technical statements must be checked independently before book publication.

\section{Shapes and Axes}

A computational tensor has a shape $(d_1,\ldots,d_k)$. Each axis needs a declared
meaning. In the book's PyTorch conventions, common examples are
\[
  X_{\mathrm{dense}}:(B,D),\qquad
  X_{\mathrm{image}}:(B,C,H,W),\qquad
  X_{\mathrm{sequence}}:(B,T,D).
\]

\section{Batched Affine Layers}

For $X\in\mathbb{R}^{B\times D_{\mathrm{in}}}$,
$W\in\mathbb{R}^{D_{\mathrm{out}}\times D_{\mathrm{in}}}$, and
$b\in\mathbb{R}^{D_{\mathrm{out}}}$,
\[
  Y=XW^\top+b.
\]
The bias broadcasts across the batch axis.

\section{Computational and Mathematical Usage}

Machine-learning software uses \emph{tensor} for a multidimensional computational
array. Such an array may store coordinates of a mathematical tensor, but the two
definitions should not be treated as interchangeable without specifying coordinate
transformation rules.

% [The upstream code, diagrams, device claims, FLOP claims, performance advice,
% channels-last examples, and autograd survey are omitted.]

\end{document}
